\chapter{Discussion et limites}
\label{chapitre:discussion}

% =====================================================================
\section{Points forts de l'approche}
% =====================================================================

\subsection{Positionnement méthodologique clair}
Le système ne prétend pas diagnostiquer cliniquement la boiterie; il produit une alerte de
boiterie probable, interprétable et exploitable comme aide à la décision. Ce positionnement
explicite réduit le risque de surpromesse scientifique et respecte une exigence récurrente
de la littérature: les systèmes de détection automatisée doivent être évalués dans leur
contexte opérationnel plutôt que comparés à des standards cliniques dont les labels sont
souvent absents ou incertains \citep{alsaaod2019automatic}.

\subsection{Architecture hybride IF + règles métier}
L'ablation montre qu'un détecteur d'anomalies seul produit des alertes impulsives
(variante~C: 100\% de détection mais aucune précision temporelle), tandis qu'un système
purement à règles devient rigide et générateur de faux positifs (variante~B: $3.2\times$
plus de fausses alertes). Le couplage retenu permet d'obtenir des épisodes plus plausibles
et des notifications plus défendables, conformément à l'approche préconisée par les travaux
récents en élevage de précision \citep{michelena2025plf_review}.

\subsection{Profondeur de validation}
Peu de projets de cette envergure dépassent la simple démonstration visuelle. Ici, la
validation combine injections manuelles
(Tableau~\ref{tab:manual_v21}), grilles robustes multi-seeds sur 9\,504~exécutions au total
(Tableau~\ref{tab:protocole_validation}), comparaison inter-versions avec 91.7\% d'accord
(Tableau~\ref{tab:versions_comparison}), benchmark performance
(Tableau~\ref{tab:performance}) et vérification SHA des artefacts. Cette chaîne rend le
projet beaucoup plus crédible qu'une simple interface appuyée sur quelques exemples choisis.

\subsection{Honnêteté de l'ablation}
L'ablation étendue apporte une discussion utile sur le choix du détecteur. La comparaison
IF/LOF (Chapitre~\ref{chapitre:validation}) révèle des profils complémentaires: IF domine
sur la précision temporelle des cas forts, tandis que LOF capte mieux les cas borderline.
IF est maintenu comme détecteur principal pour sa meilleure précision sur les cas
actionnables et sa complexité algorithmique favorable ($O(n \log n)$ contre $O(n^2)$).
Une telle transparence illustre la démarche du mémoire: documenter les alternatives testées
plutôt que les masquer.

Le Tableau~\ref{tab:forces_limites} synthétise les principales forces et limites discutées
dans les sections suivantes.

\begin{table}[H]
    \centering
    \caption{Synthèse des forces et limites du système proposé.}
    \label{tab:forces_limites}
    \begin{tabular}{p{6.5cm}p{6.5cm}}
        \toprule
        \textbf{Forces} & \textbf{Limites} \\
        \midrule
        Positionnement explicite (alerte, pas diagnostic) &
            Pas de validation clinique au sens strict \\
        Architecture hybride IF + règles (interprétable) &
            Fragilité sur les cas borderline ($\sim$42\% en validation robuste finale) \\
        Validation multi-niveaux (9\,504~exécutions) &
            Mono-ferme, mono-capteur, 41~jours \\
        Ablation honnête avec comparaison LOF &
            Validation LOF limitée au protocole v4 \\
        Scalabilité linéaire ($\approx$0.6\,s/vache) &
            IF représente 89\% du temps de calcul \\
        Traçabilité SHA-256 et reproductibilité &
            Seuils dépendants du corpus courant \\
        \bottomrule
    \end{tabular}
\end{table}

\FloatBarrier
% =====================================================================
\section{Limites}
% =====================================================================

\subsection{Absence de labels cliniques exhaustifs}
La validation est technique et méthodologique, pas clinique au sens fort. Les injections
synthétiques simulent un comportement compatible avec la boiterie, mais ne sont pas des cas
cliniques confirmés par un vétérinaire. Cette limite est inhérente au contexte du projet:
les labels cliniques synchronisés par tranche de 15~minutes n'existent pas dans le corpus
disponible. Elle est toutefois partagée par la majorité des travaux publiés dans le domaine
\citep{oleary2020accelerometers,qiao2021review}.

\subsection{Cas borderline encore fragiles}
Les scénarios limites (borderline) restent moins bien capturés que les cas forts.
Sur l'ablation élargie (Tableau~\ref{tab:ablation_by_profile}), la détection \textit{any}
pour les cas borderline atteint seulement 27\% pour IF + règles et 40\% pour LOF + règles.
Un tel constat rejoint les observations de \cite{chapinal2010measurement}, qui
rapportent que les boiteries légères (score~2) sont les plus difficiles à détecter
automatiquement, même avec des capteurs de haute précision. Il doit être assumé comme une
limite du système.

\subsection{Dépendance aux seuils}
Même si les seuils ont été explorés et justifiés via une grille de 336~configurations
(v5.2), ils restent liés au contexte expérimental du corpus courant. Un changement de
capteur, de race, de saison ou de conditions d'élevage pourrait nécessiter un
recalibrage. Ce constat est partagé par \cite{thorup2015lameness}, qui observent des
variations significatives de performance entre fermes avec les mêmes capteurs. Le gel des
seuils est une mesure de reproductibilité, pas une garantie d'universalité.

\subsection{Généralisation à confirmer}
\label{subsec:generalisation}

Les artefacts montrent une bonne stabilité interne sur le corpus présent
(28~vaches, 37\,839~lignes, 41~jours; $\sigma$ inter-seed de détection
sur la configuration finale: voir Tableau~\ref{tab:stabilite_seeds}),
mais cette stabilité ne préjuge pas de la généralisation. La littérature
rappelle que l'hétérogénéité des protocoles entre études reste un
obstacle majeur au transfert d'un détecteur d'une exploitation à l'autre
\citep{alsaaod2019automatic,oleary2020accelerometers}. Plus précisément,
cinq axes de généralisation restent à valider avant tout déploiement
clinique:

\begin{enumerate}
    \item \textit{Multi-fermes.} Le corpus provient d'une seule
    exploitation. Les profils d'activité dépendent de la conduite
    d'élevage (fréquence de traite, accès au parcours, conception des
    stalles), ce qui peut modifier la distribution des \textit{features}
    de base (\texttt{Steps}, \texttt{Motion\,Index}, transitions) et
    donc recalibrer implicitement les seuils de règle et le contour de
    l'Isolation Forest. Un plan réaliste consisterait à ré-évaluer la
    configuration figée sur au moins 3~fermes indépendantes avant
    d'énoncer des performances cibles.

    \item \textit{Multi-capteurs.} Les données proviennent exclusivement
    du podomètre IceQube. Un changement de capteur (AfiAct, Smartbow,
    CowManager, accéléromètre brut au collier) modifie la définition
    opérationnelle des variables: un \og pas \fg{} IceQube n'est pas un
    \og pas \fg{} AfiAct, et le \textit{Motion Index} n'existe pas chez tous les
    fournisseurs. Les \textit{features} dérivées (\texttt{*\_rrz},
    \texttt{*\_d1\_per\_hour}, \texttt{*\_vs\_mm7}) sont en revanche
    définies de manière \textit{capteur-agnostique} à partir de séries
    temporelles brutes, ce qui laisse une voie de portage, mais qui
    doit être testée empiriquement.

    \item \textit{Multi-saisons.} La fenêtre d'observation
    (octobre--novembre) couvre une période automnale. Les schémas
    comportementaux peuvent différer en été (stress thermique, pâturage)
    ou en hiver (stabulation prolongée, luminosité réduite). La
    normalisation robuste rolling (\texttt{*\_rrz}) absorbe une partie
    des variations lentes, mais des effets de niveau persistants
    pourraient dégrader l'AUC de 0.956 (\S\ref{subsubsec:roc_pr})
    observée sur ce corpus.

    \item \textit{Validation clinique des injections.} Les profils
    d'ablation (\textit{strong}, borderline,
    \textit{multisignal\_persistent}) sont construits sur des
    hypothèses plausibles mais non calibrées contre des cas cliniques
    réels scorés par locomotion. Le chiffre d'\textit{detection\_any}
    de 47.8~\% sur 90~cas injectés ne doit donc pas être lu comme une
    sensibilité clinique. La baseline pédométrique d'Alsaaod
    (\S\ref{subsubsec:alsaaod_baseline}) est elle aussi évaluée sur ces
    mêmes injections, ce qui rend la \textit{comparaison} valide
    (mêmes cas, même moteur de règles) sans pour autant trancher la
    question de la performance absolue en conditions réelles.

    \item \textit{Dérive temporelle.} Sur 41~jours, les dérives lentes
    (vieillissement du capteur, changement d'alimentation, variation
    parité/stade de lactation) ne peuvent pas être observées. Un suivi
    longitudinal de 6~à~12~mois permettrait d'estimer la fréquence de
    ré-étalonnage nécessaire (\S\ref{subsec:workflow_ferme} a proposé
    un scénario de dérive calibration, mais son seuil de déclenchement
    est pour l'instant fixé sans données longitudinales).
\end{enumerate}

En résumé, les performances rapportées au
Chapitre~\ref{chapitre:validation} sont des performances
\textit{internes} à un corpus méthodologique, et non des performances
attendues \textit{en champ}. Le Tableau~\ref{tab:forces_limites} cadre
ce positionnement et la section \og Perspectives \fg{} de la
conclusion en tire les priorités de recherche.

\subsection{Biais potentiels du corpus}
Le corpus de 28~vaches sur 41~jours, bien que suffisant pour une validation méthodologique,
présente des biais potentiels qui doivent être assumés:
\begin{itemize}
    \item \textit{Biais saisonnier}: les données couvrent octobre--novembre, une période
    automnale où les conditions de logement et de parcours peuvent influencer les
    comportements de manière spécifique. Les schémas comportementaux pourraient
    différer en été (pâturage, chaleur) ou en hiver (stabulation prolongée);
    \item \textit{Biais de sélection}: les 28~vaches du corpus ne constituent pas un
    échantillon représentatif au sens épidémiologique. La composition du troupeau
    (parité, stade de lactation, race) n'est pas documentée dans les données capteurs;
    \item \textit{Biais d'injection}: les profils synthétiques simulent des déviations
    comportementales plausibles mais pas cliniquement validées. L'amplitude et la forme
    des perturbations injectées reposent sur des hypothèses raisonnables mais non
    calibrées contre des cas cliniques réels.
\end{itemize}

De tels biais ne remettent pas en cause la validité de la démarche méthodologique
(validation reproductible, ablation, comparaison IF/LOF), mais ils délimitent la portée
des conclusions: les performances rapportées sont valides pour ce corpus, dans ces
conditions, et la généralisation reste à confirmer.

\FloatBarrier
% =====================================================================
\section{Menaces à la validité}
% =====================================================================

Toute étude empirique est exposée à des menaces qui peuvent limiter la portée de ses
conclusions. Cette section complète les limites identifiées ci-dessus en les organisant
selon le cadre classique de \cite{shadish2002experimental}, adapté aux études en génie
logiciel par \cite{wohlin2012experimentation} et \cite{feldt2010validity}.

\textit{Validité interne.}
Les injections synthétiques simulent des déviations comportementales plausibles mais ne
reproduisent pas la pathophysiologie réelle de la boiterie: un animal cliniquement boiteux
peut présenter des modifications plus complexes \citep{whay2003assessment}. Les profils
injectés constituent donc une approximation qui favorise potentiellement le détecteur.
Par ailleurs, les variables d'entrée sont mécaniquement corrélées (une vache couchée plus
longtemps est nécessairement moins active), ce qui amplifie le signal mais limite
l'interprétabilité de la contribution individuelle de chaque variable.
Enfin, le modèle IF est entraîné sur les 60\,\% premiers bins de chaque vache
(\texttt{baseline\_ratio}\,=\,0.6), sous l'hypothèse que cette période représente un
comportement normal. Si une vache est déjà boiteuse au début de l'observation, la
baseline est contaminée et le modèle risque de normaliser un comportement pathologique.
Une telle hypothèse est raisonnable dans un contexte de suivi continu (la boiterie clinique
est un événement qui évolue dans le temps), mais elle n'a pas été vérifiée empiriquement
sur le corpus principal faute de labels cliniques synchrones.

\textit{Validité externe.}
Les données proviennent d'une seule exploitation, d'un seul type de capteur (IceQube),
sur 41~jours en période automnale. Le contexte expérimental est typique des élevages
Holstein en stabulation libre; les résultats ne sont pas directement transférables à
d'autres races, systèmes de logement ou saisons \citep{galindo2000relationships}. La
fenêtre temporelle restreinte ne permet pas non plus de capturer les variations
comportementales à long terme.

\textit{Validité de construit.}
Le système produit une alerte de ``boiterie probable'', définie opérationnellement comme
une déviation comportementale détectée par le pipeline. Ce construit ne correspond pas à
la boiterie clinique, qui nécessite un examen physique et un score de locomotion validé
\citep{sprecher1997lameness}. La métrique IoU, empruntée à la vision par ordinateur
\citep{everingham2010pascal}, suppose que le recouvrement temporel est un indicateur
pertinent de la qualité de détection, hypothèse raisonnable pour des épisodes continus
mais potentiellement inadaptée pour des boiteries intermittentes.

\textit{Validité de conclusion.}
Les 90~cas par variante d'ablation suffisent pour détecter des différences marquées
mais limitent la capacité à identifier des effets de taille modérée. Les résultats sont
présentés principalement sous forme de moyennes agrégées; une analyse appariée
complémentaire avec IC95\% a été ajoutée pour la comparaison IF/LOF, mais le mémoire
ne propose pas de modélisation statistique hiérarchique complète ni d'intervalles pour
l'ensemble des tableaux. La sensibilité de la performance à la configuration finale,
qui constitue un point unique dans l'espace des 336~paramétrisations, fait l'objet
d'une analyse locale dédiée
(Section~\ref{subsec:sensibilite_locale}) qui montre qu'aucune perturbation à
$\pm 1$~cran ne fait basculer la décision GO/NO-GO; une analyse globale de l'ensemble
de l'espace des paramètres reste hors-périmètre.

\FloatBarrier
% =====================================================================
\section{Spécificité de la détection et diagnostic différentiel}
\label{sec:specificite_detection}
% =====================================================================

Une question naturelle face à un système de détection de boiterie probable est celle
de la \textit{spécificité étiologique}: lorsqu'une alerte est émise, comment le système
distingue-t-il une boiterie d'une autre cause possible de modification comportementale?
La présente section clarifie la portée réelle des alertes produites et documente les
mécanismes par lesquels le pipeline limite (sans les éliminer) les principaux
confondants cliniques.

\subsection{Ce que détecte réellement le pipeline}

Le système ne détecte pas la boiterie au sens clinique. Il détecte une
\textit{anomalie comportementale persistante sur des variables locomotrices}, conforme
au positionnement explicite de l'introduction (\og boiterie probable \fg{}, et non
diagnostic). Trois propriétés du pipeline orientent cette détection vers des signatures
compatibles avec une boiterie plutôt que vers des anomalies comportementales génériques:

\begin{enumerate}
    \item \textit{Ciblage des variables locomotrices.} Les cinq variables utilisées
    (\textit{Motion Index}, pas, temps couché, temps debout, transitions) sont le cœur de
    la locomotion bovine. Une pathologie qui n'affecte pas la locomotion
    (ex.: mammite subclinique sans signe systémique) ne déclenchera pas le système.
    \item \textit{Filtre de persistance ($K$~=~28~bins, soit 7~h).} Les événements
    comportementaux brefs (pic d'activité isolé, bousculade, changement ponctuel de
    stalle) sont filtrés avant de pouvoir générer une notification.
    \item \textit{Cohérence multi-signaux (mode MIX).} Une vraie boiterie affecte
    typiquement plusieurs familles de signaux de manière corrélée: moins de pas,
    plus de temps couché, moins de transitions. Une perturbation qui n'affecte qu'une
    famille isolée est filtrée par la règle de cohérence.
\end{enumerate}

\subsection{Confondants cliniques et mécanismes de mitigation}

Le Tableau~\ref{tab:confondants} recense les principales causes non pathologiques ou
pathologiques non-boiteuses susceptibles de modifier le comportement locomoteur d'une
vache laitière, et indique pour chacune le mécanisme du pipeline qui limite le risque
de fausse alerte.

\begin{table}[H]
    \centering
    \caption{Confondants potentiels et mécanismes de mitigation par le pipeline.}
    \label{tab:confondants}
    \small
    \begin{tabular}{p{3.0cm}p{3.5cm}p{7.0cm}}
        \toprule
        \textbf{Confondant} & \textbf{Risque de déclenchement}
            & \textbf{Mécanisme de mitigation} \\
        \midrule
        Œstrus (chaleurs)
            & Modéré à élevé
            & Signature d'activité \textit{augmentée} (opposée à la boiterie);
              durée typique 6--18~h, parfois filtrée par la persistance de 7~h;
              croisement possible avec détection d'œstrus existante. \\
        \addlinespace
        Mammite subclinique
            & Faible à modéré
            & Signature différente (léthargie isolée sans déséquilibre locomoteur
              cohérent); la cohérence multi-signaux filtre les anomalies
              mono-famille. \\
        \addlinespace
        Vêlage
            & Faible
            & Dates de vêlage connues de l'éleveur; exclusion triviale en
              pré-traitement. La modification comportementale est par ailleurs
              très brève à l'échelle de l'horizon de persistance. \\
        \addlinespace
        Stress thermique
            & Très faible
            & Affecte l'ensemble du troupeau simultanément; la spécificité
              individuelle du pipeline (seule la vache cible réagit,
              cf.~Tableau~\ref{tab:before_after}) filtre naturellement ces
              perturbations globales. \\
        \addlinespace
        Changement de ration, stress social
            & Très faible
            & Mêmes arguments que le stress thermique: perturbation à l'échelle
              du troupeau, filtrée par la spécificité individuelle. \\
        \addlinespace
        Problème podal non boiteux (ex.: lésion débutante)
            & Modéré
            & Peut déclencher une alerte; ce déclenchement est \textit{cliniquement
              pertinent}, car il attire l'attention sur un animal à risque avant
              l'apparition d'une boiterie franche. \\
        \bottomrule
    \end{tabular}
\end{table}

\subsection{Positionnement clinique: outil de priorisation, pas de diagnostic}

La portée des alertes doit être lue comme celle d'un \textit{test de dépistage}
(screening), et non d'un test diagnostique. Un test de dépistage est conçu pour
identifier les animaux qui justifient une vérification plus approfondie; le diagnostic
différentiel définitif reste la prérogative du vétérinaire ou d'un éleveur expérimenté,
sur la base d'un examen clinique, d'un score de locomotion \citep{sprecher1997lameness}
ou d'un examen podal. Cette division du travail est cohérente avec la pratique du
dépistage en médecine humaine: un test positif ne signe pas la maladie, il justifie
une investigation.

Concrètement, une alerte produite par le système doit être interprétée comme:
\og cette vache présente une déviation comportementale persistante sur ses variables
locomotrices; une vérification terrain est recommandée \fg{}. Les alertes émises par
le système déclenchent dans l'interface l'affichage d'un plan d'action clinique
temporisé (Section~\ref{sec:interface_streamlit}), qui oriente l'utilisateur vers
les étapes d'un examen vétérinaire structuré.

\subsection{Limites assumées}

Deux limites doivent être reconnues explicitement. D'une part, le pipeline ne permet
pas, en l'état, de distinguer algorithmiquement une boiterie d'un œstrus prolongé
sans information complémentaire: seule la signature statistique (activité
\textit{augmentée} typique de l'œstrus vs \textit{diminuée} typique de la boiterie)
offre une piste de différenciation, qui n'a pas été intégrée comme règle formelle.
D'autre part, l'absence de \textit{gold standard} clinique dans le corpus empêche
toute mesure quantitative de la spécificité étiologique du système; les arguments
présentés ci-dessus sont qualitatifs et s'appuient sur la plausibilité
mécanistique, non sur une validation empirique.

L'intégration de sources de données cliniques complémentaires (production laitière,
température corporelle, dates de vêlage, diagnostics vétérinaires) constitue un axe
d'amélioration naturel, cohérent avec les perspectives développées en
Chapitre~\ref{chapitre:conclusion}.

\FloatBarrier
% =====================================================================
\section{Mise en perspective avec la littérature}
% =====================================================================

Le système proposé se situe dans la famille des approches hybrides non-supervisées
appliquées à la détection de boiterie. À titre de comparaison, Lavrova et al.
\citep{lavrova2023lameness} obtiennent une sensibilité d'environ 77\% avec un modèle
logistique à effets mixtes sur six fermes allemandes, en combinant activité, couchage et
production laitière. Le présent projet atteint 90\% de détection sur les cas forts
(Tableau~\ref{tab:ablation_by_profile}), mais dans un contexte mono-ferme et sans labels
cliniques. La comparaison directe des performances est donc délicate, conformément à la
mise en garde d'Alsaaod et al. \citep{alsaaod2019automatic} sur l'hétérogénéité des
protocoles.

Michelena et al. \citep{michelena2025comparative} montrent que le choix de l'algorithme non
supervisé dépend du compromis visé entre sensibilité et faux positifs. Ce constat est
exactement celui que l'ablation illustre dans ce projet: IF et LOF représentent deux points
différents sur cette courbe de compromis, sans supériorité universelle de l'un sur l'autre.

La revue de Chandola et al. \citep{chandola2009anomaly} sur la détection d'anomalies
souligne que les méthodes basées sur l'isolation et celles basées sur la densité locale
ont des profils de performance complémentaires: les premières excellent sur les anomalies
globales (points isolés dans l'espace des \textit{features}), les secondes sur les anomalies
contextuelles (points normaux globalement mais déviants localement). La validation
exploratoire IF/LOF de ce projet confirme empiriquement cette distinction théorique.

Par rapport à des approches supervisées comme celles d'Alsaaod et al.
\citep{alsaaod2012electronic}, qui obtiennent des performances élevées mais sur des
cohortes labellisées cliniquement, le présent projet se distingue par sa capacité à
fonctionner en l'absence de labels, tout en maintenant une validation multi-niveaux
explicite. Cette différence de paradigme est structurelle: la détection non supervisée est
moins précise au sens clinique, mais plus déployable dans les contextes réels où les labels
cliniques sont rares ou absents \citep{qiao2021review}.

Enfin, les travaux récents de Palma et al. \citep{palma2025scoping} sur l'IA en élevage
laitier identifient la boiterie comme l'un des cas d'usage où la combinaison capteurs +
règles métier offre le meilleur rapport entre déployabilité et fiabilité. Le présent
projet s'inscrit directement dans cette famille d'approches.

Le Tableau~\ref{tab:comparaison_litterature_ch4} (Chapitre~\ref{chapitre:validation})
situe le présent projet par rapport aux travaux les plus proches. La sensibilité élevée sur
les cas forts (90--100\%) doit être relativisée: elle est mesurée sur des injections
synthétiques, pas sur des cas cliniques confirmés. En revanche, aucun des travaux cités ne
propose une validation multi-niveaux aussi étendue (9\,504~exécutions, ablation 4~variantes,
vérification d'intégrité), ce qui constitue la contribution distinctive de ce mémoire.

\FloatBarrier
% =====================================================================
\section{Comparaison avec les systèmes commerciaux}
% =====================================================================

Il est pertinent de situer le système par rapport aux solutions commerciales déployées en
élevage. Le Tableau~\ref{tab:systemes_commerciaux} compare trois systèmes représentatifs.

\begin{table}[H]
    \centering
    \small
    \caption{Comparaison du système proposé avec des solutions commerciales de suivi
    comportemental bovin.}
    \label{tab:systemes_commerciaux}
    \begin{tabular}{>{\raggedright\arraybackslash}p{2cm}
                    >{\raggedright\arraybackslash}p{2.2cm}
                    >{\raggedright\arraybackslash}p{2.0cm}
                    >{\raggedright\arraybackslash}p{2.1cm}
                    >{\raggedright\arraybackslash}p{2.2cm}}
        \toprule
        \textbf{Système} & \textbf{Capteur} & \textbf{Algorithme}
        & \textbf{Transparence} & \textbf{Validation publiée} \\
        \midrule
        CowManager
        & Boucle auriculaire
        & Propriétaire (boîte noire)
        & Faible
        & Études indépendantes limitées \\
        \addlinespace
        IceQube / IceTag
        & Patte (accéléromètre)
        & Seuils + modèles internes
        & Partielle
        & \cite{beer2016objective} \\
        \addlinespace
        Heatime / SCR
        & Collier (rumination + activité)
        & Propriétaire
        & Faible
        & \cite{stangaferro2016use} \\
        \addlinespace
        \textbf{Ce mémoire}
        & Accéléromètre
        & IF + règles\\(code ouvert)
        & \textbf{Totale}
        & \textbf{9\,504~exécutions, ablation, SHA} \\
        \bottomrule
    \end{tabular}
\end{table}

Le contraste principal est structurel: là où les systèmes commerciaux fonctionnent comme des
boîtes noires rarement évaluées de manière indépendante \citep{rutten2013invited}, le
pipeline proposé est entièrement documenté et validé quantitativement. En contrepartie, ces
systèmes bénéficient d'avantages opérationnels (intégration aux logiciels de traite,
calibration multi-fermes, support technique) que le prototype académique ne possède pas.

\FloatBarrier
% =====================================================================
\section{Implications pratiques pour l'éleveur}
% =====================================================================

En configuration finale, le système produit 0.105~notification par vache-jour
(Tableau~\ref{tab:robust_final}), soit environ 10 à 11~alertes quotidiennes
pour un troupeau de 100~vaches. Ce volume reste en dessous du seuil critique au-delà duquel
les éleveurs cessent de consulter les alertes ($\sim$40\% de faux positifs,
\citep{rutten2013invited}), et demeure compatible avec une consultation matinale avant la
première traite.

L'asymétrie des erreurs justifie l'arbitrage opérationnel retenu. Un faux positif coûte
5--10~minutes de vérification; un faux négatif se traduit par un animal souffrant non
traité à temps, avec un coût estimé entre 75 et 300~USD par cas
\citep{green2002financial,rodrigues2017lameness}. Le choix de privilégier la précision
temporelle (IF + règles) sur la sensibilité maximale (LOF seul) vise à préserver la
confiance de l'éleveur dans le système.

Le Tableau~\ref{tab:routine_eleveur} présente un scénario détaillé d'intégration dans
la journée type d'un élevage de 100~vaches (troupeau Holstein, stabulation libre).

\begin{table}[H]
    \centering
    \small
    \caption{Scénario de journée type pour un élevage de 100~vaches.}
    \label{tab:routine_eleveur}
    \begin{tabular}{lp{3.4cm}p{3.0cm}p{4.0cm}}
        \toprule
        \textbf{Moment} & \textbf{Action} & \textbf{Durée estimée}
            & \textbf{Sortie attendue} \\
        \midrule
        5h45 & Exécution automatique du pipeline & $\approx$1~min (100~vaches)
            & Fichier d'alertes horodaté \\
        \addlinespace
        6h00 & Consultation des alertes (interface ou CSV) & 5--8~min
            & Liste de 10--11 vaches prioritaires, avec score de confiance \\
        \addlinespace
        6h15 & Traite du matin: observation ciblée & 2--3~min/vache signalée
            & Score locomoteur rapide (1--5) pour chaque vache signalée \\
        \addlinespace
        7h00 & Décision vétérinaire ou traitement & Selon gravité
            & Appel vétérinaire si score~$\geq$~3; soins locaux si score~2 \\
        \addlinespace
        14h00 & Traite du soir: re-vérification rapide & 1~min/vache signalée
            & Confirmer la résolution ou escalader \\
        \addlinespace
        Vendredi & Bilan hebdomadaire & 10~min
            & Récidives, tendances troupeau, ajustement éventuel \\
        \bottomrule
    \end{tabular}
\end{table}

Un tel scénario mobilise environ 30 à 40~minutes supplémentaires par jour pour un éleveur
de 100~vaches. Appliqué à une exploitation déjà équipée de capteurs IceQube, ce qui
est le cas typique, le pipeline s'intègre sans acquisition de nouveau matériel.

% ---- SUBSECTION: Analyse coût-bénéfice ----
\subsection{Analyse coût-bénéfice}

La viabilité économique d'un système de détection précoce repose sur l'équilibre entre
les coûts évités grâce aux interventions plus rapides et les coûts opérationnels liés à
la vérification des alertes. Le Tableau~\ref{tab:cout_benefice} présente ce bilan pour
trois tailles de troupeau, sous les hypothèses conservatrices suivantes~: prévalence
25\,\%, coût moyen par cas non détecté de 200~USD, réduction de 40\,\% avec détection
précoce, 0.5~h/jour pour la vérification des alertes.

\textit{Hypothèses de calcul.}
Les coûts évités s'appuient sur les estimations de \cite{green2002financial} (perte de
production 270--574~kg/lactation) et de \cite{bruijnis2010assessing} (coûts vétérinaires
directs et réforme anticipée). Une détection précoce réduit le coût par cas de 30 à 50\%
en permettant une intervention avant l'aggravation clinique, soit une économie de
60--100~USD par cas détecté à temps \citep{rodrigues2017lameness}. Le coût opérationnel
est estimé à 30~min/vache-jour d'alerte à 20~USD/h, mais chaque notification
correspond le plus souvent à un épisode en cours et non à une nouvelle vache (le taux
de \textit{cas distincts} est significativement inférieur au nombre brut d'alertes).

\begin{table}[H]
    \centering
    \small
    \caption{Estimation indicative du coût-bénéfice du système pour différentes tailles
    de troupeau.}
    \label{tab:cout_benefice}
    \begin{tabular}{>{\centering\arraybackslash}p{1.8cm}
                    >{\centering\arraybackslash}p{1.8cm}
                    >{\centering\arraybackslash}p{2.2cm}
                    >{\centering\arraybackslash}p{2.2cm}
                    >{\centering\arraybackslash}p{1.8cm}
                    >{\centering\arraybackslash}p{2.0cm}}
        \toprule
        \textbf{Troupeau} & \textbf{Cas/an} & \textbf{Coûts évités (USD)} &
        \textbf{Coût vérif./an (USD)} & \textbf{Coût non-détectés} &
        \textbf{Solde net (USD)} \\
        \midrule
        50~vaches  & 12  & +960  & 1\,825 & $-$240  & \textbf{$-$1\,105} \\
        100~vaches & 25  & +2\,000 & 3\,650 & $-$500  & \textbf{$-$2\,150} \\
        200~vaches & 50  & +4\,000 & 7\,300 & $-$1\,000 & \textbf{$-$4\,300} \\
        \bottomrule
    \end{tabular}
\end{table}

\textit{Interprétation.}
Le solde est négatif dans ces hypothèses, ce qui reflète honnêtement la limite du
prototype actuel: les alertes sont comptées individuellement par bin de 15~minutes, alors
que l'éleveur ne vérifierait en pratique qu'une fois par jour les animaux signalés.
En appliquant le coût de vérification au nombre de \textit{vaches distinctes signalées}
plutôt qu'au nombre d'alertes brutes, l'équilibre s'améliore sensiblement. De plus,
l'analyse n'inclut pas les bénéfices indirects de la détection précoce: réduction
de la réforme involontaire, amélioration de la fertilité, bien-être animal, et
réduction de la charge mentale liée à la surveillance manuelle quotidienne.

Au final, ces chiffres soulignent deux axes prioritaires pour les travaux futurs:
(i) un regroupement intelligent des alertes par épisode et par animal pour réduire la
charge de vérification, et
(ii) une calibration individuelle par vache (voir Section~\ref{subsec:analyse_par_vache})
pour limiter les sur-alertes sur les animaux naturellement actifs. Avec ces améliorations,
un troupeau de 100~vaches pourrait dépasser le seuil de viabilité économique.

\FloatBarrier
\subsection{Architecture opérationnelle de bout-en-bout}
\label{subsec:workflow_ferme}

Au-delà de la journée type (Tableau~\ref{tab:routine_eleveur}), la viabilité
opérationnelle du système dépend de trois propriétés qui n'apparaissent ni
dans les métriques événementielles ni dans le bilan coût-bénéfice: la
\emph{latence} de l'alerte, la \emph{distribution des responsabilités}
entre l'éleveur, le vétérinaire et le technicien, et la \emph{gestion
des exceptions}. Cette sous-section explicite ces propriétés pour un
déploiement réaliste sur une ferme équipée de capteurs IceQube
connectés à un logiciel de gestion de troupeau.

\paragraph{Budget de latence.}
La latence de bout-en-bout, c'est-à-dire le temps écoulé entre le moment où une
anomalie comportementale se produit dans l'étable et le moment où
l'éleveur reçoit la notification, se décompose en quatre composantes
dont les ordres de grandeur sont résumés au
Tableau~\ref{tab:latence_budget}.

\begin{table}[H]
    \centering
    \small
    \caption{Budget de latence du système de bout-en-bout
    (100~vaches, fréquence 15~min, capteurs IceQube).}
    \label{tab:latence_budget}
    \begin{tabular}{lp{5.4cm}c}
        \toprule
        \textbf{Étape} & \textbf{Description} & \textbf{Ordre de grandeur} \\
        \midrule
        Acquisition capteur & Agrégation 15~min sur le collier
            & $\leq$15~min \\
        \addlinespace
        Téléchargement & Passerelle Wi-Fi ou cellulaire $\to$ base
            locale & 1--5~min \\
        \addlinespace
        Exécution pipeline & Features + IF + règles
            (0.6~s/vache, cf.~\S\ref{subsec:scalabilite})
            & 1--2~min \\
        \addlinespace
        Notification & Écriture CSV, push e-mail ou interface
            & $<$1~min \\
        \midrule
        \textbf{Total} & Budget de latence bout-en-bout
            & \textbf{20--25~min} \\
        \bottomrule
    \end{tabular}
\end{table}

Un tel budget est compatible avec l'usage visé (consultation matinale avant
la traite) mais n'est pas compatible avec une détection en temps réel
de chutes ou de ruptures accidentelles, pour lesquelles le délai de
15~min d'agrégation du capteur constitue un plancher incontournable.
Toute latence inférieure à 20~min nécessiterait un changement de
résolution capteur (non couvert par le matériel actuel) ou une
détection sur bins bruts sans fenêtre de confirmation, au prix d'une
explosion attendue des faux positifs.

\paragraph{Rôles et responsabilités.}
Le système ne produit qu'un signal de triage; toute action clinique
reste sous responsabilité humaine. Le
Tableau~\ref{tab:raci_workflow} précise la matrice RACI-lite
(\emph{Responsable, Approuve, Consulté, Informé}) pour les étapes
clés du processus.

\begin{table}[H]
    \centering
    \small
    \caption{Matrice de responsabilités pour le workflow de détection.}
    \label{tab:raci_workflow}
    \begin{tabular}{p{4.6cm}cccc}
        \toprule
        \textbf{Étape} & \textbf{Éleveur} & \textbf{Vétérinaire}
            & \textbf{Technicien} & \textbf{Système} \\
        \midrule
        Maintenance capteurs        & A & -- & R & I \\
        Exécution quotidienne       & I & -- & C & R \\
        Consultation alertes        & R & -- & -- & I \\
        Vérification terrain        & R & C  & -- & I \\
        Diagnostic différentiel     & C & R  & -- & -- \\
        Décision de traitement      & A & R  & -- & -- \\
        Saisie retour clinique      & R & C  & -- & I \\
        Recalibration annuelle      & C & -- & R & -- \\
        \bottomrule
    \end{tabular}
\end{table}

Un point important ressort de cette matrice: le système n'est
\emph{responsable} que de deux étapes (exécution, notification), et
n'a jamais de rôle d'approbation. Ce positionnement, cohérent avec
la catégorie ``outil de priorisation, pas de diagnostic'' défendue en
Section~\ref{sec:specificite_detection}, limite l'exposition
réglementaire du projet, puisqu'il ne relève pas d'un dispositif
médical vétérinaire au sens du règlement européen 2017/746 ou de son
équivalent canadien. L'éleveur reste l'unique responsable de la
vérification terrain, le vétérinaire du diagnostic, et l'équipe
technique de la maintenance capteurs, soit trois compétences déjà
présentes dans une exploitation équipée de capteurs.

\paragraph{Plan de contingence.}
Trois scénarios d'exception méritent un protocole explicite pour
éviter qu'une défaillance dégrade la confiance dans l'ensemble du
système:
\begin{enumerate}
    \item \textit{Panne capteur sur une ou plusieurs vaches.} Le
    pipeline détecte naturellement l'absence de données d'entrée
    (aucun bin pour une vache donnée sur une fenêtre de 24~h) et
    n'émet aucune alerte pour cette vache, plutôt qu'une fausse
    négative silencieuse. Un rapport quotidien liste les vaches
    sans signal, déclenchant une vérification matérielle par le
    technicien. L'absence d'alerte sur un animal sans capteur
    \emph{ne doit jamais être interprétée comme une absence de
    boiterie}: c'est la raison d'être du rapport parallèle.
    \item \textit{Indisponibilité du vétérinaire.} Lorsque le
    vétérinaire référent ne peut être joint, l'alerte reste active
    dans l'interface sur 72~h (équivalent d'un long week-end) et
    l'éleveur peut escalader vers un second vétérinaire via le
    contact de secours. Les alertes répétitives sur la même vache
    sont fusionnées pour éviter la fatigue d'alerte, conformément
    à l'axe (i) identifié dans l'analyse coût-bénéfice.
    \item \textit{Dérive de calibration.} La calibration
    (Section~\ref{subsubsec:calibration}) est vérifiée en continu
    par comparaison de la distribution empirique du score
    \texttt{if\_anom\_k\_max} à la distribution de référence du
    corpus d'entraînement. Un test de Kolmogorov-Smirnov glissant
    à $p < 0.01$ déclenche un drapeau de recalibration; tant que
    le drapeau est levé, les probabilités calibrées affichées à
    l'éleveur sont remplacées par le score ordinal brut et une
    mention explicite (\emph{``calibration en cours de
    vérification''}).
\end{enumerate}

\paragraph{Boucle de rétroaction clinique.}
La Figure~\ref{fig:usage_flow} (p.~\pageref{fig:usage_flow}) suggère
une boucle de retour d'expérience comme perspective d'amélioration.
Concrètement, chaque vérification terrain produit un triplet
\emph{(alerte, score locomoteur observé, décision prise)} qui,
agrégé sur 3 à 6 mois, permettrait deux raffinements: (a)~un
ré-entraînement annuel de la régression isotone de calibration sur
les \emph{vraies} étiquettes cliniques (et non plus la distinction
strong vs borderline simulée), et (b)~une calibration par vache
individuelle (\S\ref{subsec:analyse_par_vache}) informée par le
taux de vrais positifs historique de cet animal. Aucun des deux
raffinements ne modifie l'architecture: ils se branchent sur la
même régression isotone déjà validée sur corpus synthétique, avec
un apprentissage continu transparent pour l'éleveur.

\FloatBarrier
% =====================================================================
\section{Perspectives de déploiement}
% =====================================================================

Le passage d'un prototype académique à un outil opérationnel en ferme nécessite de combler
plusieurs écarts. Le benchmark de performance (0.6~s/vache) garantit une exécution
quotidienne rapide, mais l'intégration aux logiciels de gestion de troupeau, la calibration
multi-fermes et la validation clinique terrain restent à réaliser. Le
Tableau~\ref{tab:conditions_deploiement} synthétise ces conditions.

\begin{table}[H]
    \centering
    \caption{Conditions de déploiement opérationnel et statut dans le projet actuel.}
    \label{tab:conditions_deploiement}
    \begin{tabular}{lcc}
        \toprule
        \textbf{Condition} & \textbf{Statut} & \textbf{Effort estimé} \\
        \midrule
        Performance ($<$5~min/troupeau) & Satisfait & -- \\
        Reproductibilité (seeds, SHA) & Satisfait & -- \\
        Calibration mono-ferme & Satisfait & -- \\
        Intégration logiciel de gestion & Non couvert & Moyen \\
        Calibration multi-fermes & Non couvert & Important \\
        Interface éleveur simplifiée & Prototype & Moyen \\
        Validation clinique terrain & Non couvert & Important \\
        \bottomrule
    \end{tabular}
\end{table}

% ---- Retour sur la validation exploratoire IceTag 2019 ----
\subsection{Retour sur la validation exploratoire IceTag 2019}

La Section~\ref{sec:validation_mcgill} a présenté les résultats de l'application du
pipeline gelé au corpus complémentaire IceTag (30~vaches, même ferme, automne~2019).
La portabilité technique est confirmée, et les alertes sont partiellement cohérentes
avec les antécédents cliniques pour les cas stables (vaches~821, 3444, 2078). En
revanche, le cas de la vache~5879 (saine stricte mais fortement alertée) et les
effectifs trop réduits ($n$\,=\,5 vs $n$\,=\,2) empêchent de conclure à une
validation clinique.

Une telle limitation est essentiellement imputable au décalage de 8 à 9~mois entre les
scores SLS (hiver~2019) et les données capteurs (automne~2019). Une validation
clinique synchrone, dans laquelle les évaluations locomotrices seraient contemporaines
des enregistrements capteurs, constitue la prochaine étape nécessaire. Une telle
étude permettrait de quantifier la sensibilité et la spécificité réelles du système
contre un gold standard vétérinaire, et de déterminer si les alertes observées sur
la vache~5879 reflètent une boiterie acquise entre les évaluations ou une
sur-détection du pipeline. La généralisation inter-fermes resterait également à
confirmer.

\FloatBarrier
% =====================================================================
\section{Considérations d'usage et éthique}
% =====================================================================

\subsection{Responsabilité algorithmique et bien-être animal}

L'adoption d'un système de détection crée une obligation implicite envers les animaux
surveillés \citep{broom2017animal}: un faux négatif n'est pas une simple erreur
statistique, mais un animal souffrant non identifié à temps \citep{whay2003assessment}.
Une telle asymétrie éthique justifie le choix de privilégier la précision temporelle sur
la sensibilité maximale, et impose que la notification soit toujours suivie d'une
vérification terrain, idéalement articulée avec un scoring locomoteur
\citep{sprecher1997lameness,ahdb_mobility_scoring_dairy_cows,dairynz_preventing_managing_lameness_guide_2025}.
Comme le soulignent \cite{estevez2019animal}, un système qui
génère trop de bruit peut paradoxalement dégrader le bien-être animal en créant un faux
sentiment de sécurité.

\subsection{Confiance et adoption}

La confiance dans la fiabilité des alertes est souvent plus déterminante que la performance
brute \citep{bewley2015precision,venkatesh2003user}. Le phénomène de \textit{fatigue d'alerte}
\citep{rutten2013invited} constitue le principal obstacle à l'adoption durable: les
éleveurs cessent de consulter les notifications au-delà d'un seuil critique de faux
positifs. Le présent système, avec 0.105~notification par vache-jour en configuration MIX,
reste sous ce seuil. La transparence des règles métier (seuil de score, fenêtre glissante,
cooldown) constitue un levier supplémentaire de confiance, en distinguant l'approche des
systèmes commerciaux fonctionnant comme des boîtes noires
(cf. Tableau~\ref{tab:systemes_commerciaux}).

\subsection{Protection des données}

Les données capteurs révèlent indirectement les pratiques d'élevage et constituent un actif
informationnel sensible dont la propriété reste largement non résolue
\citep{wolfert2017big}. Le présent projet, en fonctionnant entièrement en local, préserve
la souveraineté de l'éleveur sur ses données. Plus largement, trois choix de conception
répondent aux principes d'IA responsable \citep{carbonell2016ethics}: traçabilité complète
(versionnage, SHA-256), documentation explicite des limites, et architecture locale sans
centralisation sur des serveurs tiers
\citep{palma2025scoping,michelena2025plf_review}.

La Figure~\ref{fig:radar_systeme} synthétise le profil du système proposé sur six dimensions
complémentaires. Les forces (sensibilité forte, maîtrise des faux positifs, scalabilité,
interprétabilité) contrastent avec la faiblesse attendue en généralisation multi-fermes
(0.30), conformément aux limites du corpus. Ce type de représentation permet de
visualiser les forces et les faiblesses du pipeline de manière globale et de comparer
différentes configurations ou approches alternatives.

\begin{figure}[H]
    \centering
    \begin{tikzpicture}[scale=1.1]
        % Axes du radar (6 axes, 60° entre chaque)
        \foreach \i in {0,...,5} {
            \pgfmathsetmacro{\angle}{90-\i*60}
            \draw[black!20] (0,0) -- (\angle:3.5);
        }
        % Grille hexagonale
        \foreach \r in {1,2,3} {
            \draw[black!10]
                (90:\r) -- (30:\r) -- (-30:\r) --
                (-90:\r) -- (-150:\r) -- (150:\r) -- cycle;
        }
        % Labels des axes
        \node[font=\scriptsize, above] at (90:3.7) {Sensibilité cas forts};
        \node[font=\scriptsize, right] at (30:3.9) {Précision temporelle};
        \node[font=\scriptsize, right] at (-30:3.9) {Maîtrise faux positifs};
        \node[font=\scriptsize, below] at (-90:3.7) {Scalabilité};
        \node[font=\scriptsize, left] at (-150:3.9) {Interprétabilité};
        \node[font=\scriptsize, left] at (150:3.9) {Généralisation};

        % Polygone du système (valeurs normalisées 0-1)
        \filldraw[ifblue!15, draw=ifblue, thick]
            (90:3.325) -- (30:2.275) -- (-30:2.975) --
            (-90:3.15) -- (-150:3.15) -- (150:1.05) -- cycle;
        % Points sur les sommets
        \fill[ifblue] (90:3.325) circle (2.5pt);
        \fill[ifblue] (30:2.275) circle (2.5pt);
        \fill[ifblue] (-30:2.975) circle (2.5pt);
        \fill[ifblue] (-90:3.15) circle (2.5pt);
        \fill[ifblue] (-150:3.15) circle (2.5pt);
        \fill[ifblue] (150:1.05) circle (2.5pt);
    \end{tikzpicture}
    \caption{Profil radar du système sur six dimensions d'évaluation.}
    \label{fig:radar_systeme}
\end{figure}

La Figure~\ref{fig:usage_flow} illustre le positionnement du système dans le flux
décisionnaire d'un élevage~: la notification est un signal de triage, la décision finale
revenant à l'opérateur, et les flèches pointillées indiquent une boucle de retour
d'expérience envisagée comme perspective d'amélioration.

\begin{figure}[H]
    \centering
    \resizebox{\textwidth}{!}{%
    \begin{tikzpicture}[
        box/.style={rectangle, draw=black!70, minimum height=1.1cm,
                     minimum width=3cm, align=center, font=\small, inner sep=12pt},
        decision/.style={rectangle, draw=black!70, minimum width=2.8cm, minimum height=1.1cm,
                         align=center, font=\small, inner sep=12pt},
        arrow/.style={-{Stealth[length=2.5mm]}, thick, black},
        darrow/.style={-{Stealth[length=2.5mm]}, thick, dashed, black!50},
        lbl/.style={font=\scriptsize, black, fill=white, inner sep=1.5pt},
    ]
        \node[box, fill=profgray!10] (capteur) {Capteurs\\[-7pt]comportementaux};
        \node[box, fill=ifblue!12, right=1.4cm of capteur] (pipeline) {Pipeline\\[-7pt]de détection};
        \node[box, fill=loforange!12, right=1.4cm of pipeline] (notif) {Notification\\[-7pt]boiterie probable};
        \node[decision, fill=rulegreen!12, right=1.4cm of notif] (decision) {Vérification\\[-7pt]terrain};
        \node[box, fill=rulegreen!15, above right=0.8cm and 1.2cm of decision]
            (action) {Intervention\\[-7pt]vétérinaire};
        \node[box, fill=profgray!8, below right=0.8cm and 1.2cm of decision]
            (noaction) {Pas d'action\\[-7pt](faux positif)};

        \draw[arrow] (capteur) -- (pipeline);
        \draw[arrow] (pipeline) -- (notif);
        \draw[arrow] (notif) -- (decision);
        \draw[arrow] (decision.north east) -- (action.west)
            node[lbl, pos=0.5, above, sloped] {confirme};
        \draw[arrow] (decision.south east) -- (noaction.west)
            node[lbl, pos=0.5, below, sloped] {infirme};

        % Boucle retour
        \draw[darrow] (noaction.south) -- ++(0,-0.8) -| (pipeline.south)
            node[lbl, pos=0.25, below] {retour d'expérience (perspective)};
    \end{tikzpicture}%
    }
    \caption{Positionnement du système dans le flux décisionnaire d'un élevage.}
    \label{fig:usage_flow}
\end{figure}

\vspace{0.5em}
\noindent
En résumé, ce chapitre a mis en perspective les forces et les limites du système, identifié les menaces
à la validité et discuté les considérations éthiques liées à son déploiement. La conclusion
qui suit synthétise les contributions principales de ce mémoire et propose des perspectives
de travail futur à court, moyen et long terme.
